\subsection{The category of SWFs}

In this section we discuss some categorical constructions and results on $\SWF_k$.

There is a clear forgetful functor from $\SWF_k$ to the category of topological spaces, $U\colon\SWF_k\to\Top$, which maps $(X,\O_X)$ to its underlying topological space $X$
and a morphism $\phi\colon(X,\O_X)\to(Y,\O_Y)$ to its underlying continuous map $\phi$.
This functor is faithful, but not full (it is easy to construct a continuous map which is not a morphism of SWFs).
Though it is a left-adjoint, as we will show.

\bdefn

    For a topological space $X$ and an arbitrary set $S$, a function $f\colon X\to S$ is {\emphcolor locally constant} if for every $p\in X$ there is a neighborhood
    $p\in U\subseteq X$ such that $f|_U$ is constant.

\edefn

Note that $f\colon X\to S$ being locally constant is the same as it being continuous when $S$ is endowed with the discrete topology.
Indeed, if $f$ is locally constant then for every $p\in f^{-1}\set x$ there is an open neighborhood of $p$ contained in $f^{-1}\set x$ by definition, so $f^{-1}\set x$ is open.
Thus $f$ is continuous to discrete $S$.
Conversely, $f$ being continuous to discrete $S$ means that $f^{-1}\set x$ is open and so for every $p\in X$ we take the neighborhood $U=f^{-1}(fp)$ to see it is locally constant.

\blemm

    For an SWF $(X,\O_X)$, $\O_X(U)$ contains all locally constant functions $U\to k$.

\elemm

\bproof

    Let $f\colon U\to k$ be locally constant, so there is an open cover $U_i$ of $U$ on which $f|_{U_i}$ is constant.
    Since $\O_X(U_i)$ is a $k$-algebra, it must contain all constant functions (since it contains the unit and all of its multiples), so $f|_{U_i}\in\O_X(U_i)$.
    By gluing, this means $f\in\O_X(U)$ as desired.
    \qed

\eproof

\bprop

    The forgetful functor $U\colon\SWF_k\to\Top$ is a left-adjoint.

\eprop

\bproof

    Forgetful functors are in general ``the most general'' object, so in our case in lieu of the previous lemma it seems that the following construction is a good contender.
    For open $U\subseteq X$ define $\O_X(U)$ to be the collection of locally constant $U\to k$.
    Thus $\O_X$ is just the sheaf of continuous maps into $k$, which is always a sheaf.

    For locality, let $f\in\O_X(U)$ then its preimage on all subsets of $k$ is open, in particular $D_U(f)$.
    And if $f$ is locally constant, so too is $1/f$ ($f$ and $1/f$ are constant on the same open subsets).
    Thus $(X,\O_X)$ is an SWF which is denoted $\ul X$.

    We claim that $U\dashv\ul\bul$.
    This amounts to constructing a natural isomorphism
    $$ \hom_\Top(U(X,\O_X),Y) \cong \hom_\SWF((X,\O_X),\ul Y) . $$
    $U$ defines an injective map from the right to the left, so we just need to show that it is surjective.
    That is, for $\phi\colon X\to Y$ continuous, $\phi$ lifts to a morphism of SWFs $(X,\O_X)\to\ul Y$.
    So for $U\subseteq Y$ open we need to show $\phi^*$ maps $\O_Y(U)\to\O_X(\phi^{-1}U)$.

    For $f\colon U\to k$ locally constant, then $f\circ\phi\colon\phi^{-1}U\to k$ is locally constant as well, as the composition of continuous maps.
    Thus $\phi^*(f)\in\O_X(\phi^{-1}U)$ by the previous lemma.
    \qed

\eproof

In particular $U$ preserves colimits.
The following is a technical definition and lemma which we will not prove in full here (see homework).

\bdefn

    Let $\c B$ be a basis for a topology, then a {\emphcolor B-SWF} over $k$ is a presheaf $\O\colon\c B^\op\to\Alg_k$ satisfying the following axioms:
    \benum
        \item $\O$ is a subpresheaf of the function sheaf, $\fun(\bul,k)$.
        \item {\emphcolor B-sheaf axiom}: for any basic set $B\in\c B$ and cover of basic sets $B=\bigcup_iB_i$, if $f\colon B\to k$ is a function such that $f|_{B_i}\in\O(B_i)$
        for all $i$, then $f\in\O(B)$.
        \item {\emphcolor locality}: for any $B\in\c B$ and $f\in\O(B)$, $D_B(f)$ is open (i.e.\ the union of basic sets) and for every basic cover $B_i$ of $D_B(f)$, we have that
        $1/f|_{B_i}\in\O(B_i)$.
    \eenum

\edefn

\bprop

    For every B-SWF $(\c B,\O)$ where $\c B$ is a basis of the topology on $X$, there is a unique SWF $(X,\O_X)$ extending it, meaning $\O_X(B)=\O(B)$ for all $B\in\c B$.

\eprop

\bproof

    The construction is as follows: for every $U\subseteq X$ open and $\set{B_i}\subseteq\c B$ covering $U$, define
    $$ \O_X(U) = \set{f\colon U\to k}[f|_{B_i}\in\O(B_i)\hbox{ for all $i$}] . $$
    Then one can show without much difficulty that $\O_X(U)$ is independent on the choice of basic covering.
    In particular this means that $\O_X(B)=\O(B)$ for all $B\in\c B$ by choosing the covering $\set B$.

    It is then easy to show that this is a sheaf, and locality follows easily from locality of $(\c B,\O)$.
    \qed

\eproof

Now we will show that $\SWF_k$ is cocomplete.
Since $U$ preserves colimits, we already know what the underlying topological spaces of the colimits are.

\blemm

    Let $\pi\colon X\to Y$ be a quotient map, and suppose there is an SWF structure $(X,\O_X)$ on $X$.
    Then there is a unique SWF structure $(Y,\O_Y)$ given by the following universal property: $\phi\colon(Y,\O_Y)\to(Z,\O_Z)$ is a morphism of SWFs iff $\phi\circ\pi\colon(X,\O_X)\to(Z,\O_Z)$
    is a morphism of SWFs.

\elemm

\bproof

    This is a universal property and thus a solution is unique up to unique isomorphism.
    Define
    $$ \O_Y(U) = \set{f\colon U\to k}[\pi^*(f)\in\O_X(\pi^{-1}(U)] . $$
    It is not hard to show that this satisfies the universal property.
    \qed


\eproof

\bthrm

    The category $\SWF_k$ is cocomplete.

\ethrm

\bproof

    We just need to show that $\SWF_k$ has coproducts and coequalizers.
    For coproducts, let $\set{X_i}_i$ be a collection of SWFs and let $X=\coprod_iX_i$.
    It has a basis of open subsets given by the union of all open subsets of $X_i$.
    Then define $\O_X$ on $X$ as a B-SWF where for $U\subseteq X_i$ open we define
    $$ \O_X(U) = \O_{X_i}(U) . $$
    It is easy to see that this is a B-SWF, and thus defines an SWF on $X$.
    It is also not hard to see that this has the universal property of a coproduct.

    For coequalizers, let $\phi,\psi\colon X\tto Y$ be parallel morphisms of SWFs.
    The underlying space of their coequalizer will be $Y/{\sim}$ where $\sim$ is the equivalence relation generated by $\phi(x)\sim\psi(x)$ for all $x\in X$.
    This is a quotient space, so we just give it the quotient space SWF structure from the previous lemma.
    This clearly satisfies the universal property of coequalizers as a direct result from the universal property of quotients.
    \qed

\eproof

\bnote

    While the colimit of SWFs is an SWF, the colimit of algebraic varieties need not be a variety!

\enote

\bcoro

    The global sections functor $\Gamma$ contravariantly preserves colimits.

\ecoro

\bproof

    Since $\Gamma$ is contravariant, preserving colimits means mapping colimiting cones to limit cones.
    So it is sufficient to show that
    $$ \Gamma\parens{\coprod_i(X_i,\O_{X_i})} \cong \prod_i\Gamma(X_i,\O_{X_i}),\qquad \Gamma(\coeq(\phi,\psi)) \cong \eq(\Gamma(\phi),\Gamma(\psi)) , $$
    in a way which maps the colimit cone to the limit cone, which we leave as an exercise.

    For coproducts, by the above construction we see that for $X=\coprod_iX_i$,
    $$ \Gamma\parens{\coprod_i(X_i,\O_{X_i})} = \set{f\colon X\to k}[f|_{X_i}\in\O_{X_i}(X_i)] , $$
    so the map $\Gamma(X)\to\prod_i\Gamma(X_i)$ given by $f\mapsto(f|_{X_i})$ is a natural isomorphism.

    For a quotient map $\pi\colon X\to Y$, we have
    $$ \Gamma(Y) = \set{f\colon Y\to k}[f\circ\pi\in\O_X(X)] , $$
    meaning $\pi^*\colon\Gamma(Y)\to\Gamma(X)$ is an injection (by $\pi$'s surjectivity), meaning $\Gamma(Y)\cong\im\pi^*$.
    So letting $Z=\coeq(\phi,\psi)$ we have
    $$ \Gamma(\coeq(\phi,\psi)) = \set{f\colon Z\to k}[f\circ\pi\in\O_Y(Y)] \cong \im\pi^* . $$
    We know that
    $$ \eq(\Gamma\phi,\Gamma\psi) = \ker(\phi^*-\psi^*) . $$
    So we want to show that
    $$ \im\pi^* = \ker(\phi^*-\psi^*) . $$
    Firstly, we have $\pi\circ\phi=\pi\circ\psi$ by definition so $\phi^*\circ\pi^*=\psi^*\circ\pi^*$ and thus $(\phi^*-\psi^*)\circ\pi^*=0$.
    So we have left-to-right inclusion.
    Now suppose that $f\circ\phi=f\circ\psi$ for $f\in\Gamma(Y)$, then by the equivalence on $Z$ this means $f$ lifts to a map $g\colon Z\to k$, i.e.\ $f=g\circ\pi$.
    So by definition $g\in\Gamma(Z)$ since $f\in\Gamma(Y)$, and this shows that $f\in\im\pi^*$ as desired.
    \qed

\eproof

Let us consider a specific type of colimit.
Recall that a pushout is a colimit diagram of the form

\centerline{\drawdiagram{
    $A$&$B$\cr
    $C$&$X$
}{
    \diagarrow{from={1,1}, to={1,2}}
    \diagarrow{from={1,1}, to={2,1}}
    \diagarrow{from={2,1}, to={2,2}, dashed}
    \diagarrow{from={1,2}, to={2,2}, dashed}
}}

For sets, and by extension topological spaces, $X$ is the quotient of $B\amalg C$ obtained by identifying the images of an element from $A$ in $B$ and $C$.
The topology is given by the quotient topology.
This kind of colimit is called a {\emphcolor pushout}, it is dual to the {\emphcolor pullback} (also called the {\emphcolor fiber product}).
It is the colimit of a functor on the diagram $\bul\ot\bul\to\bul$.

If our category is $\Top$, then the pushout can be viewed geometrically as gluing together $B$ and $C$ along their ``intersection'' $A$.

When we are dealing with topological spaces though (and in our case SWFs), often times we would like to glue together a family of spaces and not just two.
We can generalize the pushout to account for this.

\bdefn

    Let $I$ be an indexing set, then the {\emphcolor gluing category} over $I$ is the category $\glue_I$ whose objects are $I\sqcup\binom I2$ (i.e.\ $I$ and the set of all two-element
    subsets of $I$).
    For every $\set{i,j}\in\binom I2$, $\glue_I$ has unique morphisms from $\set{i,j}$ to $i$ and to $j$.

\edefn

For example, if $I={1,2,3}$ then $\glue_I$ looks like

\bigskip
\centerline{\drawdiagram{
    &$\set{1,3}$\cr
    $\set{1,2}$&&$\set{2,3}$\cr
    $1$&$2$&$3$\cr
}{
    \diagarrow{from={1,2}, to={3,1}, curve=-2.5cm, dest orient={left, top}}
    \diagarrow{from={1,2}, to={3,3}, curve=2.5cm, dest orient={right, top}}
    \diagarrow{from={2,1}, to={3,1}}
    \diagarrow{from={2,1}, to={3,2}}
    \diagarrow{from={2,3}, to={3,2}}
    \diagarrow{from={2,3}, to={3,3}}
}}

And when $I$ has two elements then $\glue_I$ is just the pushout category $\bul\ot\bul\to\bul$.

A {\emphcolor gluing datum} is a functor whose domain is a gluing category, i.e.\ a functor $X\colon\glue_I\to\C$.
Equivalently it is given by a collection of objects $X_i\in\C$ and $X_{ij}\in\C$ for $i,j\in\C$ ($X_{ij}$ is unordered), as well as morphisms $f_{ij}\colon X_{ij}\to X_i$ and $f_{ji}\colon X_{ij}\to X_j$.
A {\emphcolor gluing} is the colimit of a gluing datum.
Geometrically, for a gluing datum $X\colon\glue_I\to\C$, we interpret $X_i$ as the spaces we want to glue, and $X_{ij}$ as the common subspace over which we want to glue.

We know that $\SWF_k$ is cocomplete so it has all gluings.
Our main concern regarding gluings is that we can view an SWF as the gluing of an open cover of induced SWFs.

Denote the inclusion (canonical) map $X_i\to\coprod_iX_i$ by $\inc_i$.

\blemm

    For a gluing datum $X\colon\glue_I\to\C$, its gluing is the coequalizer of the diagram (where $X_{ij}=X_{ji}=X_{\set{i,j}}$)
    $$ \coprod_{i\neq j}X_{ij} \xtto\phi[\psi] \coprod_iX_i \to \colim X . $$
    Where $\phi$ is the map whose component on $X_{ij}$ is $f_{ij}$ composed with inclusion $X_i\to\coprod_iX_i$.
    Similarly $\psi$ is the map whose component on $X_{ij}$ is $f_{ji}$ composed with inclusion $X_j\to\coprod_iX_i$.

\elemm

\bproof

    Cones under $X$ are equivalent (classically) to cones under the coequalizer diagram
    $$ \coprod_{\set{i,j}\to i}X_{ij} \xtto\alpha[\beta] \coprod_iX_i\amalg\coprod_{\set{i,j}}X_{ij}, $$
    where $\alpha$'s component on $X_{ij}$ for $\set{i,j}\to i$ is $f_{ij}$ composed with inclusion into the large coproduct (on the right).
    And $\beta$'s component on $X_{ij}$ is just the inclusion into the large coproduct.

    Explicitly this means that
    $$ \alpha\circ\inc_{ij\to i} = \inc_i\circ f_{ij},\qquad \beta\circ\inc_{ij\to i} = \inc_{\set{i,j}},\qquad \phi\circ\inc_{ij} = \inc_i\circ f_{ij},\qquad \psi\circ\inc_{ij} = \inc_j\circ f_{ji} . $$

    Cones under this are maps $f$ from the large coproduct to an object $Y$ such that $f\circ\alpha=f\circ\beta$.
    We claim that restricting $f$ to $\coprod_iX_i$ gives a cone under the diagram in the lemma.
    We have that
    $$ f\circ\alpha\circ\inc_{ij\to i} = f\circ\beta\circ\inc_{ij\to i} $$
    which is equivalent to
    $$ f\circ\inc_i\circ f_{ij} = f\circ\inc_{\set{i,j}} . $$
    Thus we see that this is also equal to $f\circ\inc_j\circ f_{ji}$.

    A cone under the diagram in the lemma is a map $f$ such that
    $$ f\circ\phi\circ\inc_{ij} = f\circ\psi\circ\inc_{ij} \iff f\circ\inc_i\circ f_{ij} = f\circ\inc_j\circ f_{ji} $$
    for all $i,j$.
    So we have shown that a cone under the classical coequalizer diagram restricts to one of the one in the lemma.

    Conversely if we have $f\circ\inc_i\circ f_{ij}=f\circ\inc_j\circ f_{ji}$ for $f\colon\coprod_iX_i\to Y$, then this uniquely extends to a cone over the classical diagram.
    This is because we must have
    $$ f\circ\inc_i\circ f_{ij} = f\circ\inc_{\set{i,j}} . $$
    Thus we can just define $f$'s component on $X_{ij}$ to be $f\circ\inc_i\circ f_{ij}$, and by the identity for $f\circ\phi=f\circ\psi$ this is well-defined (symmetric in $i$ and $j$).

    So there is an isomorphism between the cones under the classical and the diagram in the lemma, so they define the same coequalizers.
    \qed

\eproof

\blemm

    Let $X\colon\glue_I\to\Top$ be a gluing datum.
    Then its gluing is
    $$ \colim X = \coprod_{i\in I}X_i/{\sim} , $$
    where $\sim$ is the equivalence induced by $f_{ij}(x)\sim f_{ji}(x)$ for all $x\in X_{ij}$.

\elemm

\bproof

    Consider the coequalizer diagram in the previous diagram
    $$ \coprod_{i\neq j}X_{ij} \xtto\phi[\psi] \coprod_iX_i \to \colim X . $$
    The coequalizer is just the quotient of $\coprod_iX_i$ by identifying the images of $\phi$ and $\psi$ of the same element.
    $\phi$ on $X_{ij}$ is just $f_{ij}$ and $\psi$ is $f_{ji}$, so this is just the equivalence $f_{ij}(x)\sim f_{ji}(x)$ for $x\in X_{ij}$ as desired.
    \qed

\eproof

Since we know what the coproduct of SWFs and the quotient of an SWF is, we can then determine the SWF obtained by gluing.

Let $X\colon\glue_I\to\SWF_k$ be a gluing datum.
Write $X_i=(X_i,\O_i)$, $X_{ij}=(X_{ij},\O_{ij})$, and $\phi_{ij}=X(\set{i,j}\to i)$ for the data.
Let us write $\hat X=\coprod_iX_i$.
Then its gluing is the SWF $(X,\O_X)$ where $X=\hat X/{\sim}$ where $\phi_{ij}(x_{ij})\sim\phi_{ji}(x_{ij})$.
Let us denote by $\pi\colon\hat X\to X$ the projection.

For $U\subseteq X$ open, then $\pi^{-1}U\subseteq\hat X$ is open and thus has the form $\pi^{-1}U=\bigcup_iU_i$ where $U_i=X_i\cap\pi^{-1}U$ are open.
Then by our construction of quotient and coproduct SWFs,
$$ \O_X(U) = \set{f\colon U\to k}[f\circ\pi\in\O_{\hat X}(\pi^{-1}U)] = \set{f\colon U\to k}[(f\circ\pi)|_{U_i}\in\O_i(U_i)\hbox{ for all $i$}] . $$

We can now prove our goal regarding gluing.

\bthrm

    Let $(X,\O_X)$ be an SWF and $X=\bigcup_iU_i$ an open cover.
    Then $(X,\O_X)$ is isomorphic to the gluing of $(U_i,\O_X|_{U_i})$ along their intersections.

\ethrm

\bproof

    This is a rather lengthy proof.
    It is entirely technical, and just ensures that everything behaves as it should.

    Let us write $\O_i$ for $\O_X|_{U_i}$ and $\O_{ij}$ for $\O_X|_{U_i\cap U_j}$.
    Define $(Y,\O_Y)$ to be the gluing of the $U_i$ along their intersections.
    Let $\hat Y=\coprod_iU_i$, so as a topological space $Y=\hat Y/{\sim}$ where we identify elements of $U_i,U_j$ if they lie in the intersection.
    For an element $x\in U_i$ let $x_i\in\hat Y$ denote its copy in $\hat Y$.
    Then our equivalence is induced by $x_i\sim x_j$ for $x\in U_i\cap U_j$.
    Furthermore for $x\in\hat Y$ let $\pi(x)=[x]\in Y$ denote its equivalence class, and for $x\in U_i\subseteq X$ we define $[x]=[x_i]$.

    There is a clear map $\hat\phi\colon\hat Y\to X$ which maps an element of $\hat Y$ to its corresponding element in $X$, i.e.\ $x_i\mapsto x$.
    This is a continuous map, since for $U\subseteq X$ open we have $\hat\phi^{-1}(U)=\coprod_i(U_i\cap U)$ open.
    Furthermore it respects the equivalence $\sim$, since if $x\in U_i\cap U_j$ then $x_i$ and $x_j$ are both mapped to $x$.
    So by the universal property of quotients, this defines a map $\phi\colon Y\to X$ given by $\phi[x]=x$ (so we have shown that such a map is well-defined and continuous).
    Clearly $\phi$ has an inverse given by $x\mapsto[x]$, so it is bijective.

    Now we claim that $\hat\phi$ is an open map.
    This is clear since it maps basic open sets $U\subseteq U_i$ to $U\subseteq X$ which are open (since $U_i$ is open in $X$).
    So we have a commutative diagram

    \bigskip
    \centerline{\drawdiagram{
        $\hat Y$&&$X$\cr
        &$Y$
    }{
        \diagarrow{from={1,1}, to={1,3}, text=$\hat\phi$, y distance=.25cm}
        \diagarrow{from={1,1}, to={2,2}, text=$\pi$, x distance=-.25cm}
        \diagarrow{from={2,2}, to={1,3}, text=$\phi$, y distance=-.25cm}
    }}

    Every open set of $Y$ is of the form $\pi(\hat U)$ for open $\hat U\subseteq\hat Y$ (since $U\subseteq Y$ is open iff $\pi^{-1}U\subseteq\hat Y$ is open, and $U=\pi(\pi^{-1}U)$ by surjectivity).
    And we see $\phi(\pi(\hat U))=\hat\phi(\hat U)$ is open.
    So $\phi$ is a bijective open continuous map, thus a homeomorphism.

    Now we need to show that for all $V\subseteq X$ open, $\phi^*$ defines an isomorphism $\O_X(V)\to\O_Y(\phi^{-1}V)$.
    As we saw with B-SWFs, it is sufficient to restrict our attention to basic open sets.
    We take as a basis for $X$ the open subsets of $U_i$, as $i$ ranges.
    These define a basis of open subsets of $\hat Y$, and thus projecting them via $\pi$ defines a basis of $Y$.

    For $V\subseteq U_i$ we have that restricting $\hat\phi$ to $V\subseteq\hat Y$ gives the identity $V\to V\subseteq X$.
    That is
    $$ V = \hat\phi(V) = \phi\circ\pi(V) \implies \phi^{-1}(V) = \pi(V) . $$
    So if we show that $\phi^*$ is an isomorphism $\O_X(V)\to\O_Y(\phi^{-1}V)=\O_Y(\pi(V))$ then we will have finished, as we have exhausted both the basic open sets of $X$ as well as $Y$'s.

    By above we know for $U\subseteq Y$ open we have
    $$ \O_Y(U) = \set{f\colon U\to k}[(f\circ\pi)|_{\hat U_i}\in\O_i(\hat U_i)\hbox{ for all $i$}] , $$
    where $\hat U_i=\pi^{-1}(U)\cap U_i$.
    In our case, for $V\subseteq U_i$ basic we are concerned with $U=\pi(V)$.
    Then $\hat U_j=\pi^{-1}(\pi(V))\cap U_j$ which are all $x\in U_j$ such that $x$ is identified with an element of $V\subseteq U_i$.
    So we see $\hat U_j=V\cap U_j$ and so
    $$ \O_Y(\pi(V)) = \set{f\colon\pi(V)\to k}[(f\circ\pi)|_{V\cap U_j}\in\O_j(V\cap U_j)\hbox{ for all $j$}] . $$
    Since $V\cap U_j$ form an open cover of $V$, this is just equal to
    $$ \O_Y(\pi(V)) = \set{f\colon\pi(V)\to k}[(f\circ\pi)|_V\in\O_X(V)] . $$

    We know that on $V$, $\hat\phi$ is just the identity so $\phi\circ\pi=\id$ and since $\phi$ is bijective, $\phi$ and $\pi$ are inverses.
    So for $f\in\O_X(V)$ we have $\phi^*(f)\circ\pi=f\circ\phi\circ\pi=f$ thus $\phi^*(f)\in\O_Y(\pi(V))$.
    So $\phi$ is indeed a morphism of SWFs.

    And $\phi^*$ is bijective, as we can map $f\in\O_Y(\pi(V))$ to $f\circ\pi\in\O_X(V)$.
    Then $\phi^*(f\circ\pi)=f\circ\pi\circ\phi=f$ as they are inverses on $V$, and $\phi^*(f)\circ\pi=f\circ\phi\circ\pi=f$.
    Thus $\phi$ is an isomorphism as desired.
    \qed

\eproof

Our main interest in the above theorem is due to algebraic varieties.
By definition, algebraic varieties are SWFs $X$ with an affine cover $U_i$.
This theorem says we can view $X$ as a certain colimit of the $U_i$ with $U_i\cap U_j$.
In particular we can use the fact that $\Gamma$ preserves colimits to compute the global sections of $X$ as a certain limit of $\Gamma(U_i)$ and $\Gamma(U_i\cap U_j)$.

And if $Y\subseteq X$ is a subspace, then by the above theorem $Y$ is the gluing of $Y\cap U_i$ along their intersections $Y\cap U_i\cap U_j$.
Thus global sections of $Y$, which are just sections of $Y$, are a certain limit of $\Gamma(Y\cap U_i)$ and $\Gamma(Y\cap U_i\cap U_j)$.
In particular for $V\subseteq X$ open, $\O_X(U)$ is a certain limit of $\O_X(U\cap U_i)$ and $\O_X(U\cap U_i\cap U_j)$.

So if we know $U_i$ and the intersections $U_i\cap U_j$ well, we can deduce properties of $X$.

We now consider another functorial construction on $\SWF_k$.

\bdefn

    Let $(X,\O_X)$ be an SWF and $p\in X$ a point.
    Define its {\emphcolor stalk} at $p$ to be
    $$ \O_{X,p} = \colim_{x\in U}\O_X(U) , $$
    where the colimit is taken of the functor $\O_X$ restricted to the subposet $\N(p)\subseteq\Op(X)$ of open neighborhoods of $p$.
    Explicitly, this is
    $$ \O_{X,p} = \left.\coprod_{x\in U}\O_X(U)\middle/{\sim}\right. , $$
    where $\sim$ is the equivalence identifying $f\in\O_X(U)$ and $g\in\O_X(V)$ if there is an open subset $x\in W\subseteq U,V$ such that $f|_W=g|_W$ (i.e.\ $f$ and $g$ are equal
    locally around $p$).
    Such equivalence classes are called {\emphcolor germs}.

\edefn

Since the colimit of $k$-algebras is a $k$-algebra, the stalk of an SWF is a $k$-algebra.

We sometimes write $(f,U)$ to mean $f\in\O_X(U)$ or to mean its germ in $\O_{X,p}$.
Notice that for a germ $(f,U)$ of $p$, the value $f(p)$ makes sense, as all functions in the same germ must have the same value at $p$.

Consider the category of {\emphcolor pointed spaces of functions}: $\SWF_k^*$, whose objects are SWFs with a distinguished point: $(X,\O_X,p)$ for $p\in X$,
and a morphism of pointed SWFs is a morphism of SWFs which preserve the distinguished points: $\phi\colon(X,\O_X,p)\to(Y,\O_Y,q)$ is a morphism of SWFs such that $\phi(p)=q$.

Then taking stalks ({\emphcolor stalkification}) is functorial $(\SWF_k^*)^\op\to\Alg_k$.
This is because a morphism of pointed SWFs $\phi\colon(X,\O_X,p)\to(Y,\O_Y,q)$ induces a morphism of functors $\phi^*\colon\O_Y\to\O_X$ as well as a functor $\phi^{-1}\colon\N(q)\to\N(p)$ via
mapping $q\in V$ to $\phi^{-1}V$.
Thus it forms a morphism of functors $\O_Y|_{\N(q)}\to\O_X|_{\N(p)}$ and since taking colimits is itself functorial, we get a morphism $\O_{Y,q}\to\O_{X,p}$.

Explicitly $\phi$ induces the map $\phi^*\colon\O_{Y,q}\to\O_{X,p}$ on germs by mapping $(f,U)\in\O_{Y,q}$ to $(\phi^*(f),\phi^{-1}U)\in\O_{X,p}$.

Further notice that if $p\in U\subseteq X$ is a neighborhood of $p$ then the stalks of $\O_X|_U$ and $\O_X$ at $p$ are isomorphic.
This can be shown by the germ definition, or by noting that the stalk is a filtered colimit, and $\N(p,U)^\op$ is a cofinal subset of $\N(p,X)^\op$, and filtered colimits over the restriction
of a functor to a cofinal subcategory is isomorphic to the colimit over the whole subcategory.

\blemm

    Let $Y$ be an affine variety and $P\in Y$.
    Then under the isomorphism $Y\cong\specm(k[Y])$, $P$ corresponds to the maximal ideal $\m_P\leq k[Y]$ given by
    $$ \m_P = \set{f\in k[Y]}[f(P)=0] . $$

\elemm

\bproof

    As $Y$ is affine, it is of the form $Z_{\bA^n}(I)$ for some ideal $I\leq k[\bA^n]$, $\I(Y)=I$.
    The isomorphism $Z_{\bA^n}(I)\cong\specm(k[Y])$ is given by mapping $P=(a_1,\dots,a_n)$ to $\m_P=(x_1-a_1,\dots,x_n-a_n)$, quotiented by $I$.
    Clearly $\m_P\subseteq\set{f\in k[Y]}[f(P)=0]$, and the right-hand side is a proper ideal so there is equality.
    \qed

\eproof

\bcoro

    Let $(X,\O_X)$ be an algebraic variety, $P\in X$, and $P\in U$ an affine neighborhood.
    Then $\O_{X,P}\cong k[U]_{\m_P}$ where $\m_P\leq k[U]$ is the ideal corresponding to $P$ as in the above lemma.

\ecoro

\bproof

    As explained, taking the filtered colimit over a cofinal subsets doesn't affect the result, and so we can assume $X=U$ is affine.
    Furthermore, $\set{D(f)}_{f\in\O_X(X)}$ is a cofinal subset because it is a basis.
    Thus we can take the colimit over $D(f)$ where $f(P)\neq0$, i.e.
    $$ \O_{X,P} \cong \colim_{f\in k[X],f(P)\neq0}\O_X(D(f)) \cong \colim_{f(P)\neq0}k[X]_f . $$
    For $f(P)\neq0$ we have $f\notin\m_P$ and so it is in the localizing set of $k[X]_{\m_p}$.
    In particular, we have natural maps $k[X]_f\to k[X]_{\m_p}$ for all $f$.
    We now show that $k[X]_{\m_p}$ satisfies the universal property of the colimit.

    Suppose that we have maps $k[X]_f\to A$ for all $f$ commuting with restrictions.
    In particular these each give maps $k[X]\to A$ which map $f$ to a unit.
    Since the maps $k[X]_f\to A$ commute with restrictions, we see that this is a single map $k[X]\to A$ mapping every $f(P)\neq0$, so $f\notin\m_P$, to a unit.
    By the universal property of localization this induces a map $k[X]_{\m_P}\to A$.
    Thus $k[X]_{\m_P}$ is the colimit, as desired.
    \qed

\eproof

\bnote

    Notice that this shows that $\O_{X,P}$ is a local ring for an algebraic variety $X$.
    This is true in general for an SWF: let $p\in X$ arbitrary then let $\m\leq\O_{X,p}$ be the set of all germs $(U,f)$ such that $f(p)=0$.
    This is clearly an ideal of $\O_{X,p}$, and it is proper as it doesn't contain the unit.
    Now for $(U,f)\notin\m$ we know that $p\in D_U(f)$ and thus $(D_U(f),1/f)$ is a germ in $\O_{X,p}$ and a multiplicative inverse for $(U,f)$.
    So $O_{X,p}-\m=\O_{X,p}^\times$ so $(\O_{X,p},\m)$ is a local ring.

    This lends to a generalization of SWFs.
    As discussed in the definition of SWFs, the sheaf condition generalizes naturally to give a structure called a sheaf.
    These are presheaves valued in the category of rings (or $R$-algebras) satisfying the sheaf condition (a certain diagram is an equalizer diagram).

    But this does not explain how to generalize the locality condition.
    Utilizing stalks, we can.
    Call a pair $(X,\O_X)$ where $X$ is a topological space and $\O_X$ a sheaf valued in the category of rings (resp.\ $\Alg_R$) a {\emphcolor ringed space} (resp.\ {\emphcolor $R$-ringed space}).
    This generalizes the sheaf condition.

    Now a ($R$-) ringed space $(X,\O_X)$ is called a {\emphcolor locally ringed space} if all of its stalks $\O_{X,p}$ are local rings (resp.\ local $R$-algebras).
    This generalizes the locality condition.

\enote


